% Source: Wald GR, physical PDF p. 269
%         (printed p. 256).
% Coverage: Figure 10.3, the notion of extrinsic curvature.
% Figures/tables/footnotes: Figure 10.3; no tables or footnotes.
% Status: source-reviewed and compile-clean.
% Uncertainties: none.
\begingroup
\renewcommand{\thefigure}{10.3}
\begin{figure}[htbp]
\centering
\begin{tikzpicture}[scale=1.0,>=Latex,font=\small]
  \draw[thick] (-3.1,-.7) .. controls (-1.5,.55) and (1.0,.38) .. (3.1,.75);
  \coordinate (Q) at (-1.35,.05);
  \coordinate (P) at (.45,.42);
  \fill (Q) circle (1.5pt) node[below] {\(q\)};
  \fill (P) circle (1.5pt) node[below right] {\(p\)};
  \draw[->,thick] (Q) -- +(.18,1.15) node[above] {\(n^a\)};
  \draw[->,thick] (P) -- +(.10,1.18) node[above] {\(n^a\)};
  \draw[->,thick,dashed] (P) -- +(.52,1.05);
  \node[right] at (2.42,.62) {\(\Sigma\)};
\end{tikzpicture}
\caption{A spacetime diagram illustrating the notion of the extrinsic curvature
of a hypersurface \(\Sigma\). The dashed arrow at \(p\) represents the parallel
transport of the normal vector, \(n^a\), at \(q\) along a geodesic connecting
\(q\) to \(p\). The failure of this vector to coincide with \(n^a\) at \(p\)
corresponds intuitively to the bending of \(\Sigma\) in the spacetime in which it
is embedded. The formula \(K_{ab}=h_a{}^c\nabla_cn_b\) shows that \(K_{ab}\)
directly measures this failure of the two vectors at \(p\) to coincide for \(q\)
near \(p\).}
\label{fig:10.3}
\end{figure}
\endgroup
